% =====================================================================
\chapter{Bevezetés}
\label{ch:bevezetes}
% =====================================================================

Ez a fejezet a problémát, a \Mimir{} jelenlegi állapotát, a kutatási kérdéseket, a hipotéziseket és a
hozzájárulásokat mutatja be.

\section{Motiváció}

A feleletválasztós kérdés (\emph{multiple-choice question}, MCQ) gépi úton javítható és nagy létszámú
kurzusokon is skálázható, előállítása azonban drága: egy jó kérdés törzse önmagában is érthető, pontosan egy
válaszlehetősége helyes a tananyag szerint, a helytelen válaszok (\emph{disztraktorok}) pedig hihetők, de
igazolhatóan hibásak~\cite{haladyna2002review,gierl2017distractors}. A nagy nyelvi modellek (\emph{large language
model}, LLM) óta a legelterjedtebb megoldás, hogy az oktató egy felhőben futó csevegő asszisztenstől kér
kérdéseket a jegyzetéből. A modellek képesek kérdéseket és disztraktorokat fogalmazni~\cite{kurdi2020systematic},
de hajlamosak a forrás által nem alátámasztott állításokra~\cite{ji2023hallucination}. Ezt a munkafolyamatot a
továbbiakban \emph{közvetlen megközelítésnek} nevezzük.

\section{Problémafelvetés}

A közvetlen megközelítés kényelmes, de az oktatási alkalmazás szempontjából négy gyengesége van.

\begin{description}[leftmargin=1.1cm,labelwidth=0.9cm,style=sameline,format=\normalfont\bfseries]
\item[(P1)] \textbf{A helyesség nincs ellenőrizve.} Semmi sem garantálja, hogy a megjelölt helyes választ (a
\emph{kulcsot}) a dokumentum alátámasztja, és hogy a disztraktorok valóban hamisak; az önellenőrzés ezt
csökkentheti, de nem szünteti meg~\cite{huang2024selfcorrect}.
\item[(P2)] \textbf{A szerkezet és a lefedettség nincs kontrollálva.} A kért darabszám, típus és nehézség
csak kérés a promptban, és a hosszú bemenetet a modellek egyenetlenül használják~\cite{liu2024lost}, így a
kérdések a dokumentum néhány részére tömörülhetnek.
\item[(P3)] \textbf{Adatvédelem.} A tananyag személyes adatot is tartalmazhat; a felhőszolgáltató az
adatkezelés címzettjévé válik~\cite{gdpr2016}, és az intézmény elveszíti a kontrollt a tartalom felett.
\item[(P4)] \textbf{Függés a nagy, távoli modellektől.} A teljes dokumentum egyetlen promptban hosszú
kontextust és nagy modellt igényel, a saját, fogyasztói hardveren futtatható kis modellek kontextusa és
memóriája viszont korlátos.
\end{description}

A P1 és a P2 a \emph{módszer} hiányából fakad: ha a vizsga szerkezete, a kérdések forrása és az ellenőrzés nem
önálló lépések, sem mérni, sem kikényszeríteni nem lehet őket. A P3 és a P4 egymást erősíti: a helyi futtatás
megoldaná az adatvédelmi problémát, de a kis modellek éppen a hosszú, egylépéses feladatban a leggyengébbek.
Kiinduló gondolatunk ezért az, hogy rövid, egyenként ellenőrizhető lépésekre bontva egy fogyasztói grafikus
kártyán (GPU) futó kis modell is megbízható vizsgát állíthat elő.

\section{A rendszer állapota és a dolgozat célja}

A \Mimir{} fejlesztési fázisban lévő rendszer. Éles szolgáltatásai ma a naiv visszakereséses megoldást
(\emph{naiv RAG}, \arm{E0}) futtatják: a modell a feladatleírással visszakeresett három dokumentumrészletet
(darabot) kapja egy promptban. A lépésekre bontott tervezett folyamat a saját kiértékelő keretrendszerünkben
készült el, ahol minden lépése külön mérhető. A dolgozat célja \emph{mérésen alapuló döntés} arról, melyik
megoldás kerüljön a rendszerbe -- a legjobb, a leggyorsabb és a legrobusztusabb, külön a 8 és a 12\,GB-os
GPU-szinten. A mérések a kiinduló gondolatot pontosították; a dolgozat fő állítása:

\begin{quote}
\itshape Fogyasztói GPU-n futó kis modellekkel a vizsgakérdések minőségét elsősorban a választott modell és a
visszakeresés lekérdezése határozza meg; a feladat lépésekre bontása gyenge modellnél nélkülözhetetlen, erős
modellnél pedig az ellenőrizhetőséget adja. Ezt egy saját, nyílt kiértékelő keretrendszerrel, 30 dokumentumon
és három modellen mértük, és ebből döntöttük el, mi kerül a rendszerbe.
\end{quote}

\section{Kutatási kérdések és hipotézisek}

\begin{description}[leftmargin=1.3cm,labelwidth=1.1cm,style=sameline,format=\normalfont\bfseries]
\item[RQ1] A strukturált generálás -- tervezés, kérdésenkénti forráshoz kötés és ellenőrzés -- érvényesebb és
jobban lefedő vizsgát ad-e, mint ha a modellnek a dokumentumot (vagy annak visszakeresett részleteit)
egyetlen promptban adjuk át?
\item[RQ2] Mennyivel járulnak hozzá a minőséghez az egyes komponensek (tervező, ellenőrző, fogalomgráf,
hibrid keresés, darabolás), és mennyibe kerülnek futási időben és modellhívásban?
\item[RQ3] Mennyire közelíti meg egy fogyasztói GPU-n (8--12\,GB) futó, 4 bitre kvantált kis modell egy nagy,
szerveren futó modell teljesítményét, és hogyan függ az eredmény a választott helyi modelltől?
\end{description}

A hipotéziseket \az{\ref{sec:hipotezis-ertekeles}}.~szakasz értékeli; a konfigurációk (\arm{E0}, \arm{E1},
\dots) a kísérlet egy-egy beállítását jelölik (\ref{sec:karok}.~szakasz).

\begin{hipotezis}[Tervezés]\label{h:terv}
Azonos generátormodell mellett a vizsgaterv (\arm{E1}) a naiv visszakereséses alapvonalhoz (\arm{E0}:
szemantikus darabolás, a teljes vizsgához három visszakeresett darab) képest növeli a kérdések
alátámasztottságát és a vizsga lefedettségét, a formai megfelelés romlása nélkül.
\end{hipotezis}

\begin{hipotezis}[Ellenőrzés]\label{h:verif}
Az ellenőrző lánc (\arm{E2}) a tervezéshez (\arm{E1}) képest tovább növeli az alátámasztottságot és a vak
megoldhatóságot; az ellenőrző ítélete pozitívan egyezik a független bíró ítéletével ($\kappa>0$).
\end{hipotezis}

\begin{hipotezis}[Lefedettség]\label{h:lefed}
A teljes dokumentumot egyetlen promptban kapó alapvonal lefedettsége -- modellmérettől függetlenül --
alacsonyabb, mint a tervezett változatoké.
\end{hipotezis}

\begin{hipotezis}[Helyi vs.\ szerver]\label{h:helyi}
A helyi kis modell a tervezett és ellenőrzött folyamatban a szervermodell naiv változatának minőségi
indexét megközelíti (különbség $<0{,}1$), miközben fogyasztói GPU-n fut.
\end{hipotezis}

\begin{hipotezis}[Mérés]\label{h:meres}
A generátoroktól eltérő modellcsaládba tartozó nyelvi modell (bíró) ítéletei az oktatói értékelésekkel
legalább elfogadható megbízhatósággal egyeznek (Krippendorff-$\alpha\ge0{,}667$).
\end{hipotezis}

\section{Hozzájárulások}

\begin{enumerate}[label=(\arabic*),itemsep=0pt]
\item \textbf{Generálási módszer:} vizsgaterv, kérdéshelyenkénti keresés, forráshivatkozás, ellenőrző lánc
(\ref{ch:modszer}.~fejezet).
\item \textbf{Helyi architektúra és megfelelés:} helyi modell, törlődő vektortár, GDPR és MI-rendelet
(\ref{ch:modszer}.~fejezet).
\item \textbf{Saját, nyílt kiértékelő keretrendszer:} 14 egytényezős konfiguráció, független bíró, statisztika
(\ref{ch:kiertekeles}.~fejezet).
\item \textbf{A módszerek kiértékelése:} a naiv RAG törékenységének oka és a komponensek hatása
(\ref{ch:eredmenyek}.~fejezet).
\item \textbf{A modellek kiértékelése:} három helyi modell azonos dokumentumokon, két GPU-szint
(\ref{ch:modellek}.~fejezet).
\item \textbf{Mérési tanulságok:} szivárgás, csendes tartalékágak, $\kappa\approx0$, hiányzó bírói ítéletek
(\ref{ch:eredmenyek}.~fejezet).
\item \textbf{Gyakorlati eredmény:} a döntés arról, mi kerül a rendszerbe, GPU-szintenként
(\ref{ch:diszkusszio}.~fejezet).
\end{enumerate}

\section{A dolgozat felépítése}

\Az{\ref{ch:szakirodalom}}.~fejezet a hazai és nemzetközi irodalmat és a jogi keretet tekinti át, és
megfogalmazza a kutatási rést. \Az{\ref{ch:modszer}}.~fejezet a rendszert, a tervezett módszert és a megfelelés
helyzetét, \az{\ref{ch:kiertekeles}}.~fejezet a kiértékelési módszertant mutatja be. \Az{\ref{ch:eredmenyek}}.~fejezet
a módszer lépéseinek hatását, \az{\ref{ch:modellek}}.~fejezet a modelleket értékeli;
\az{\ref{ch:diszkusszio}}.~fejezet az eredményeket az irodalommal összeveti, és meghozza a döntést,
\az{\ref{ch:osszefoglalas}}.~fejezet pedig összefoglal. A függelék a módszer és az eredmények részleteit
(\ref{app-modszer}. és \ref{app-eredmeny-reszletek}.~függelék), a promptokat, a rubrikát, a reprodukálás
lépéseit és az OTDK kötelező elemeit tartalmazza.
